\providecommand{\Res}[2]{\qty[round-mode=figures,round-precision=3]{\fpeval{#1}}{#2}}
\providecommand{\Val}[1]{\num[round-mode=figures,round-precision=4]{\fpeval{#1}}}
\providecommand{\ResSci}[2]{\qty[round-mode=figures,round-precision=3,exponent-mode=scientific]{\fpeval{#1}}{#2}}
\providecommand{\Valt}[1]{\expanded{\noexpand\num[round-mode=figures,round-precision=4,exponent-mode=threshold,exponent-thresholds=-3\string:5]}{\fpeval{#1}}}
\providecommand{\NumT}[1]{\expanded{\noexpand\num[exponent-mode=threshold,exponent-thresholds=-3\string:9]}{#1}}
\providecommand{\QtyT}[2]{\expanded{\noexpand\qty[exponent-mode=threshold,exponent-thresholds=-3\string:9]}{#1}{#2}}
\subsubsection*{Question a) Courant initial}
\begin{enumerate}[label=\arabic*)]
    \item \textbf{Variables connues et recherchées :} $\Delta V_0 = \qty{\v{V0}}{V}$, $R = \QtyT{\vsi{R}}{\ohm}$ ; variable recherchée : $I_0$
    \item \textbf{Équations pertinentes :} à $t = 0$, toute la tension du condensateur est appliquée à la résistance :
    \[ \Delta V_0 = R I_0 \]
    \item \textbf{Isolation de la variable recherchée :}
    \[ I_0 = \frac{\Delta V_0}{R} \]
    \item \textbf{Résolution numérique :}
    \[ I_0 = \frac{\num{\v{V0}}}{\NumT{\vsi{R}}} = \ResSci{\v{V0}/\vsi{R}}{A} \]
    \item \textbf{Réponse :} $I_0 = \mathbf{\Res{1000*\v{V0}/\vsi{R}}{mA}}$.
\end{enumerate}

\subsubsection*{Question b) Temps pour atteindre \qty{\v{p}}{\percent} de la tension initiale}
\begin{enumerate}[label=\arabic*)]
    \item \textbf{Variables connues et recherchées :} $R = \QtyT{\vsi{R}}{\ohm}$, $C = \QtyT{\vsi{C}}{F}$, $\Delta V/\Delta V_0 = \NumT{\vsi{p}}$ ; variable recherchée : $t$
    \item \textbf{Équations pertinentes :}
    \[ \Delta V(t) = \Delta V_0\, e^{-t/(RC)} \]
    \item \textbf{Isolation de la variable recherchée :}
    \[ t = -RC \ln\left(\frac{\Delta V}{\Delta V_0}\right) \]
    \item \textbf{Résolution numérique :}
    \[ \tau = RC = \Res{\vsi{R}*\vsi{C}}{s} \qquad t = -\Valt{\vsi{R}*\vsi{C}} \cdot \ln(\NumT{\vsi{p}}) = \Res{-\vsi{R}*\vsi{C}*ln(\vsi{p})}{s} \]
    \item \textbf{Réponse :} $t = \mathbf{\Res{-\vsi{R}*\vsi{C}*ln(\vsi{p})}{s}}$.
\end{enumerate}

\subsubsection*{Question c) Énergie dissipée}
\begin{enumerate}[label=\arabic*)]
    \item \textbf{Variables connues et recherchées :} $C = \QtyT{\vsi{C}}{F}$, $\Delta V_0 = \qty{\v{V0}}{V}$, $\Delta V = \Res{\v{V0}*\vsi{p}}{V}$ ; variable recherchée : $E_{\text{diss}}$
    \item \textbf{Équations pertinentes :} l'énergie perdue par le condensateur est dissipée dans la résistance :
    \[ U_C = \tfrac{1}{2} C (\Delta V)^2 \qquad E_{\text{diss}} = U_{C,0} - U_C \]
    \item \textbf{Isolation de la variable recherchée :}
    \[ E_{\text{diss}} = \tfrac{1}{2} C \left[(\Delta V_0)^2 - (\Delta V)^2\right] \]
    \item \textbf{Résolution numérique :}
    \[ E_{\text{diss}} = \tfrac{1}{2} \cdot \NumT{\vsi{C}} \left[\num{\v{V0}}^2 - \left(\Valt{\v{V0}*\vsi{p}}\right)^2\right] = \ResSci{0.5*\vsi{C}*\v{V0}^2*(1-\vsi{p}^2)}{J} \]
    \item \textbf{Réponse :} $E_{\text{diss}} = \mathbf{\Res{1000*0.5*\vsi{C}*\v{V0}^2*(1-\vsi{p}^2)}{mJ}}$.
\end{enumerate}
